% 01_introduction.tex -- Introduction for the JCTC revision (ct-2026-01838j).
% Replaces the submitted Introduction and Related work. Every number is a
% \fs... macro from paper/macros.tex; \fsTODO{...} marks a number that will
% come from J10/J11/J12 and must be defined before the build.

\section{Introduction}
\label{sec:intro}

Machine-learned interatomic potentials are fitted to energies and forces from
electronic-structure calculations. Each label costs orders of magnitude more
than one evaluation of the fitted model. Which structures receive a label
therefore matters as much as any choice of architecture. The selection rule in
common use is query by committee~\citep{seung1992query,krogh1995neural}: train
several models, evaluate the spread of their predictions on every candidate
structure, and label the structures on which the spread is largest. The rule
entered atomistic modelling with neural-network
potentials~\citep{artrith2012high,gastegger2017machine,smith2018less,zhang2019active,schran2020committee}
and with linear and Gaussian-process
models~\citep{podryabinkin2017active,vandermause2020flare}, and it remains the
default in current active-learning
workflows~\citep{kulichenko2023uncertainty,chepkasov2026active}. For a
potential used in molecular dynamics the relevant spread is over forces, because
forces drive the dynamics and because force errors are larger and more local
than energy errors~\citep{christensen2020role}. Obtaining the force spread
exactly costs one reverse-mode pass through each committee member.

Committee disagreement is one of several uncertainty estimates available for
these models, and a substantial literature evaluates every such estimate
against the true error on held-out reference data; \citet{dai2025review}
review that literature. Ensembles are the
incumbent~\citep{lakshminarayanan2017simple}. Repeated comparisons of ensemble
spread with the actual error show that ensembles rank errors usefully but
understate the error magnitude, so the spread must be recalibrated before it
can be read as an error
bar~\citep{imbalzano2021uncertainty,kahle2022quality,busk2021calibrated,tan2023single,carrete2023deep}.
The same literature has settled how a calibrated estimate is recognised: by
reliability diagrams, by the variance of the standardised error, and by ranking
metrics such as the area under the receiver operating characteristic curve
(AUROC)~\citep{tran2020methods,pernot2022long,levi2022evaluating}. Conformal
prediction adds a distribution-free coverage guarantee on top of any such
score~\citep{hu2022conformal}. Single-model alternatives avoid the cost of
training and running several models: Gaussian-process and kernel
variances~\citep{musil2019fast,vandermause2020flare}, distances in feature
space~\citep{janet2019quantitative,zhu2023fast}, information-theoretic
scores~\citep{schwalbekoda2025information}, and estimators built into the
network
itself~\citep{wollschlager2023uncertainty,bigi2024rigidity,heid2024spatially}.
We state plainly, correcting the submitted version of this paper, that
comparing an uncertainty estimate with the true error is standard practice.
The distinction we draw is narrower. Work on cheap approximations of an
ensemble statistic validates the approximation against the exact ensemble, as
it should, because the approximation is meant to reproduce the ensemble. Two
links are therefore involved: the link from the ensemble to the error, and the
link from the approximation to the ensemble. A cheap approximation is only as
useful as the weaker of the two links, and this paper measures both links on
the same systems.

Deep ensembles are expensive to train and to run. Shallow ensembles reduce
the cost by sharing one network body, which we call the trunk, between $M$
output heads~\citep{kellner2024shallow,schafer2026shallow,vinchurkar2025gnn};
related multi-head constructions for graph networks and for committee-based
materials screening are given by \citet{carrete2023deep} and
\citet{beck2025committee}. The heads differ by their random initialisation, by
the subset of the training data each head sees, or by a loss term that pushes
the heads apart. Such a multi-head committee amortises the trunk: one forward
pass yields all $M$ energies, so the spread of the head energies costs no more
than a single model. What a multi-head committee does not amortise is the
force. Each head's force is the gradient of that head's energy, so obtaining
all $M$ force vectors exactly still costs $M$ backward passes through the
trunk. The alternative is a projection of the per-head forces onto a chosen
direction in head space, which costs one backward pass per direction
probed~\citep{beck2025committee}. Multi-head committees are also less diverse
than independently trained committees, because the heads share every feature
of the trunk, and they are correspondingly more
overconfident~\citep{kellner2024shallow,beck2025committee}. \citet{beck2025committee}
call an ensemble of separately trained models a naive committee; we call such an
ensemble an independently trained committee.

A second issue concerns how per-coordinate disagreement is reduced to one
number per structure. The score in common use sums the per-coordinate head
variance over all $3N$ coordinates of a structure with $N$ atoms. A sum over
coordinates grows with $N$ whether or not the model is uncertain about any one
of them. That this score depends on size is known. \citet{smith2018less} divide
the ensemble spread by $\sqrt{N}$ before comparing structures;
\citet{kellner2024shallow} show that the uncertainty of a total energy grows
between $\sqrt{N}$ and $N$ depending on how the per-atom errors are correlated;
and \citet{lortaraprasert2026size} show that a per-atom uncertainty averaged
over atoms and the same uncertainty taken as a maximum over atoms behave
differently as the structure grows. Every benchmark on which the summed score
has been recommended, however, uses one molecule at one fixed size. At fixed
size an extensive score and an intensive score differ by a constant, and the
dependence on $N$ cannot show. We show below that on a pool of variable size
the summed score selects the largest structures rather than the most uncertain
ones, and that a size normalisation restores parity with a per-coordinate
maximum rather than superiority over it.

The setting in which committee uncertainty is now most often proposed is a
foundation model: a potential pretrained on a large materials database and
then fine-tuned, or equipped with committee readouts, for a new
task~\citep{deng2023chgnet,batatia2023foundation,bilbrey2025foundation}.
Candidate pools in this setting are large, span many elements, and contain
structures of very different sizes, which is exactly where the size dependence
above bites. Such pools also carry a trap. The MACE-MP trunk used here was
pretrained on the MPtraj trajectories of \citet{deng2023chgnet}, so scoring an
MPtraj pool tests the committee heads out of sample but tests the trunk on its
own pretraining data. The submitted version of this paper did not state this.
The revision keeps the MPtraj scan, labels that scan as in-distribution for the
trunk, and adds a scan of the OMat24 validation
subsets~\citep{barrosoluque2024omat24}, which post-date the pretraining of the
trunk and which \citet{beck2025committee} used as an out-of-distribution test;
the standard benchmarks for universal potentials provide the
context~\citep{riebesell2025matbench}. MPtraj also contains many relaxed
crystals whose reference forces are close to zero, which favours any score
built from force magnitude; the materials results are therefore reported with
and without those structures.

This paper asks two questions, in this order. First: does committee
disagreement identify the structures on which the model's forces are most
wrong, when the error is measured against the reference forces? We call this
property the error-ranking quality of the committee. Second: how faithfully
does a cheap screening rule reproduce the decisions of the exact committee,
and at what cost? A screening rule spends a small number of backward passes on
a cheap score and skips the exact computation when that score, multiplied by a
calibration constant, falls below a threshold. The order of the two questions
matters. The second question takes the exact committee as its ground truth, so
the answer to the second question is worth having only once the first question
has been answered. Answering the first question needs no approximation at all:
it needs the per-head forces, which we cache, and the reference forces, which
the test sets contain. Answering it also needs baselines that cost nothing. We
use two zero-cost baselines: the norm of the predicted mean force, which is a
by-product of the forward pass, and the spread of the head energies, which
costs one forward pass and no backward pass. A committee statistic that costs
$M$ backward passes must beat both baselines before it is worth computing.

The findings of the paper, in plain words, are as follows. Items 1 to 3 were
in the submitted version; items 4 to 6 are new in this revision.
\begin{enumerate}
\item At fixed system size, exact committee disagreement ranks reference-force
error with AUROC $\fsFactorALo$--$\fsFactorAHi$ across six molecular test sets,
and the summed score is the best of the three structure-level scores, winning
$\fsGlobalWins$ of $\fsGlobalCells$ comparisons (Figure~\ref{fig:factora}). On
a pool of $\fsMPtrajN$ MPtraj crystals of variable size the summed score selects
structures whose median atom count is $\fsMPtrajExtensivity\times$ that of the
per-coordinate maximum, and the ordering of the scores reverses
(Figure~\ref{fig:extensivity}, Table~\ref{tab:mptraj}).
\item A deterministic estimator that spends every backward pass on the leading
directions of the head space skips $\fsLOSkipLoPct$--$\fsLOSkipHiPct$\% of
exact committee evaluations without changing a decision, and beats a randomised
low-rank sketch with the same number of backward passes on $\fsLOBeatsCV$ of
$\fsLOSystems$ systems (Figure~\ref{fig:gate}). A finite-sample law with no
free parameters predicts why the randomised sketch fails.
\item On the MPtraj pool the predicted force norm, which costs nothing, reaches
AUROC $\fsMPtrajFree$ against $\fsMPtrajGlobal$ for the exact committee, and
adding the committee to the force norm raises the AUROC by
$\fsCompMarginal$, which is not significant. Force errors grow with force
magnitude, so this is not a surprising result; what follows from it is that the
predicted force norm is the baseline every committee statistic must be scored
against. On fixed-composition molecular sets the committee beats that baseline
significantly on $\fsCompFixedSig$ of $\fsCompFixedN$ systems.
\item Every committee is now characterised by its energy and force errors and
by its calibration (Table~\ref{tab:accuracy}, Figure~\ref{fig:calibration}).
Per-atom disagreement correlates with per-atom error, with a Spearman
coefficient of \fsTODO{Spearman range of per-atom disagreement against
per-atom force error on the molecular sets}, but every shared-trunk committee
is overconfident: the variance of the standardised error is
\fsTODO{range of z-score variance for the shared-trunk committees} against
\fsTODO{z-score variance of the independently trained committees} for
independently trained committees of the same architecture.
\item On the OMat24 subsets, which are out of distribution for the trunk, the
exact committee reaches AUROC \fsTODO{OMat24 AUROC of the exact committee}
against \fsTODO{OMat24 AUROC of the predicted force norm} for the predicted
force norm, with all four foundation committees (Table~\ref{tab:ood}).
\item One round of active learning with retraining on 3BPA, in which each
selection rule chooses the same number of structures from the same pool,
changes the retrained model's force error on the unselected pool by
\fsTODO{J11 force-error change per rule relative to random selection}
relative to random selection (Table~\ref{tab:al}). Identifying high-error
structures and improving a retrained model are distinct outcomes, and the paper
reports them separately.
\end{enumerate}
The scope is one architecture family, MACE~\citep{batatia2022mace}, and
committees of $M=8$ members; smaller committees are studied by enumerating
sub-committees of the eight, and no larger committee was trained. The paper is
an evaluation of a selection rule, not a new potential, and claims are limited
to the systems and committee constructions listed in the Methods.

Every number in this paper is derived programmatically from a committed record,
and none is transcribed by hand. The protocol of each experiment was registered
and bound by a hash to a record written before that experiment produced any
output, so no hypothesis was adjusted after its outcome was seen. The
availability section (Section~\ref{sec:avail}) describes the archive, the
inputs that are not ours, and the commands that regenerate every quantity.
